\section{Experimental Results}
\label{sec:results}

\subsection{Completeness of the Two Runs}

Both the structural SCM and path ablation configurations successfully
processed all \(250\) questions. There were no missing predictions, no IDs
outside the benchmark, and no pipeline errors. Prediction coverage and the
success rate both reached \(100\%\).

The results in this section use only the following two paired artifact sets:

\begin{verbatim}
pipeline_predictions_structural_scm_v41_step5_v51.json
pipeline_predictions_path_ablation_v41_step5_v51.json
\end{verbatim}

Both prediction sets were generated by the
\texttt{2.1-query-aware-answer} runner and evaluated by Step 7 using
\texttt{2.0-current-pipeline-schema}. Older artifacts without the
\texttt{step5\_v51} suffix are excluded from the main results.

\subsection{Overall Results}

Table~\ref{tab:results} presents the principal metrics for the structural
SCM configuration. Because the paired path ablation achieved identical
values on all quality metrics, the two configurations are compared
separately in Table~\ref{tab:paired-ablation}.

\begin{table}[t]
    \centering
    \caption{Overall results of the structural SCM pipeline on \(250\)
    questions.}
    \label{tab:results}
    \begin{tabular}{lr}
        \hline
        \textbf{Metric} & \textbf{Result} \\
        \hline
        Prediction coverage & \(100.00\%\) \\
        Success rate & \(100.00\%\) \\
        Rule Recall@5 & \(69.93\%\) \\
        Event Recall@5 & \(74.66\%\) \\
        Top-1 Exact Path & \(35.20\%\) \\
        Oracle Exact Path & \(71.20\%\) \\
        Verification Accuracy & \(92.00\%\) \\
        Verification Macro-F1 & \(97.73\%\) \\
        Verification Balanced Accuracy & \(95.65\%\) \\
        Answer Token F1 & \(68.69\%\) \\
        Answer ROUGE-L F1 & \(67.67\%\) \\
        Citation F1 & \(62.86\%\) \\
        \hline
    \end{tabular}
\end{table}

The results show that retrieval achieves relatively high recall at the rule
and event levels. Verification accuracy is higher than the corresponding
path-retrieval and answer-generation scores. The substantial gap between
Top-1 Exact Path and Oracle Exact Path indicates that path ranking is a
major bottleneck in the pipeline.

\subsection{Retrieval Results}

Table~\ref{tab:retrieval-results} presents recall, precision, and hit rate at
the principal values of \(k\).

\begin{table}[t]
    \centering
    \caption{Retrieval results by \(k\).}
    \label{tab:retrieval-results}
    \begin{tabular}{lrrr}
        \hline
        \textbf{Metric} & \(\mathbf{k=1}\) & \(\mathbf{k=3}\) &
        \(\mathbf{k=5}\) \\
        \hline
        Rule Recall@\(k\) &
        \(33.12\%\) & \(65.11\%\) & \(69.93\%\) \\
        Rule Precision@\(k\) &
        \(68.40\%\) & \(44.93\%\) & \(29.04\%\) \\
        Rule Hit@\(k\) &
        \(68.40\%\) & \(85.60\%\) & \(88.40\%\) \\
        \hline
        Event Recall@\(k\) &
        \(23.37\%\) & \(66.83\%\) & \(74.66\%\) \\
        Event Precision@\(k\) &
        \(71.20\%\) & \(68.00\%\) & \(45.76\%\) \\
        Event Hit@\(k\) &
        \(71.20\%\) & \(88.00\%\) & \(89.20\%\) \\
        \hline
    \end{tabular}
\end{table}

At \(k=1\), precision for both rules and events is relatively high, whereas
recall remains low because a question may require multiple gold items along
a causal path. Increasing \(k\) to \(5\) substantially improves recall,
while precision decreases because the retriever introduces more candidates
that do not belong to the gold set.

Additional retrieval metrics are as follows:

\begin{itemize}
    \item Rule Macro Recall: \(74.41\%\);
    \item Rule Macro Precision: \(12.73\%\);
    \item Rule Macro F1: \(21.63\%\);
    \item Rule MRR: \(77.23\%\);
    \item Rule MAP: \(61.41\%\);
    \item Event Macro Recall: \(76.78\%\);
    \item Event Macro Precision: \(29.50\%\);
    \item Event Macro F1: \(42.52\%\);
    \item Event MRR: \(79.72\%\).
\end{itemize}

Overall rule precision is lower than event precision because the predictions
store not only rules on the primary path but also direct rule hits and
additional candidate rules. This result indicates that the retriever
prioritizes coverage over producing a minimal rule set.

\subsection{Causal Path Results}

Table~\ref{tab:path-results} presents the evaluation results for the primary
path and the oracle path.

\begin{table}[t]
    \centering
    \caption{Causal path evaluation results.}
    \label{tab:path-results}
    \begin{tabular}{lr}
        \hline
        \textbf{Metric} & \textbf{Result} \\
        \hline
        Top-1 Exact Path Match & \(35.20\%\) \\
        Top-1 Hop Accuracy & \(53.60\%\) \\
        Top-1 Path Length Accuracy & \(77.60\%\) \\
        Top-1 Edge Precision & \(54.40\%\) \\
        Top-1 Edge Recall & \(54.00\%\) \\
        Top-1 Edge F1 & \(54.13\%\) \\
        Top-1 Event F1 & \(61.49\%\) \\
        Top-1 Final Event Accuracy & \(49.60\%\) \\
        Top-1 Path Rule F1 & \(47.14\%\) \\
        \hline
        Oracle Exact Path Match & \(71.20\%\) \\
        Oracle Edge F1 & \(75.33\%\) \\
        Oracle Event F1 & \(76.59\%\) \\
        Oracle Final Event Accuracy & \(73.60\%\) \\
        Oracle Path Rule F1 & \(75.74\%\) \\
        \hline
        Path Ranking Error Rate & \(36.00\%\) \\
        Path Retrieval Error Rate & \(28.80\%\) \\
        \hline
    \end{tabular}
\end{table}

Top-1 Exact Path reaches only \(35.20\%\), whereas Oracle Exact Path reaches
\(71.20\%\). The absolute difference is:

\[
71.20\% - 35.20\% = 36.00
\text{ percentage points}.
\]

This finding shows that, for \(36\%\) of the samples, a candidate path that
matched the gold path more closely was present in the retrieval output but
was not selected as the primary path. In addition, \(28.80\%\) of the
samples did not contain the complete gold path in the candidate set,
reflecting an error at the retrieval stage rather than at the ranking stage.

Path Length Accuracy reaches \(77.60\%\), which is higher than Exact Path
Match. Thus, the retriever often identifies the correct number of hops but
may select the wrong seed event, mediator, outcome, or specific rule.

\subsection{Verification Results}

The query-aware verifier achieves an Accuracy of \(92.00\%\), a Macro-F1 of
\(97.73\%\), and a Balanced Accuracy of \(95.65\%\). The confusion matrix is
presented in Table~\ref{tab:verification-confusion}.

\begin{table}[t]
    \centering
    \caption{Confusion matrix for verification decisions.}
    \label{tab:verification-confusion}
    \begin{tabular}{lrrr}
        \hline
        \textbf{Gold} &
        \textbf{SUPPORTED} &
        \textbf{REJECT} &
        \textbf{UNCERTAIN} \\
        \hline
        \texttt{SUPPORTED} & \(210\) & \(0\) & \(20\) \\
        \texttt{REJECT\_DIRECT\_CLAIM} & \(0\) & \(20\) & \(0\) \\
        \texttt{UNCERTAIN} & \(0\) & \(0\) & \(0\) \\
        \hline
    \end{tabular}
\end{table}

All \(20\) direct-edge negative claims are correctly classified as
\texttt{REJECT\_DIRECT\_CLAIM}. Among the \(230\) gold
\texttt{SUPPORTED} samples, the system correctly classifies \(210\) and
assigns \(20\) to \texttt{UNCERTAIN}. No \texttt{SUPPORTED} case is
misclassified as \texttt{REJECT\_DIRECT\_CLAIM}, or vice versa.

Recall for the two classes with gold support is:

\begin{itemize}
    \item \texttt{SUPPORTED}: \(91.30\%\);
    \item \texttt{REJECT\_DIRECT\_CLAIM}: \(100.00\%\).
\end{itemize}

Precision for both classes reaches \(100\%\). This result indicates that the
verifier tends to abstain by assigning the \texttt{UNCERTAIN} label rather
than issuing an incorrect affirmative decision. However, the benchmark has
no gold samples in the \texttt{UNCERTAIN} class, so the accuracy of the
abstention mechanism cannot yet be evaluated comprehensively.

\subsection{Answer Results}

The answer-generation metrics are presented in
Table~\ref{tab:answer-results}.

\begin{table}[t]
    \centering
    \caption{Answer-generation results.}
    \label{tab:answer-results}
    \begin{tabular}{lr}
        \hline
        \textbf{Metric} & \textbf{Result} \\
        \hline
        Token Precision & \(65.34\%\) \\
        Token Recall & \(75.40\%\) \\
        Token F1 & \(68.69\%\) \\
        ROUGE-L F1 & \(67.67\%\) \\
        Normalized Exact Match & \(0.00\%\) \\
        Answer Present & \(100.00\%\) \\
        Mean number of prediction tokens & \(48.18\) \\
        Mean number of gold-answer tokens & \(45.07\) \\
        \hline
    \end{tabular}
\end{table}

Token Recall is higher than Token Precision, indicating that the answers
usually cover most of the gold content but include an additional causal
chain, explanatory material, or citations. A Normalized Exact Match of
\(0\%\) does not imply that every answer is incorrect, because predicted
answers generally add explanatory structure and \([E_i]\) markers that are
absent from the gold answers.

Token F1 and ROUGE-L F1 are similar, at \(68.69\%\) and \(67.67\%\),
respectively. This similarity suggests that the covered tokens generally
retain an ordering relatively close to that of the gold answer.

\subsection{Citation Results}

Citation results are presented in Table~\ref{tab:citation-results}.

\begin{table}[t]
    \centering
    \caption{Citation evaluation results.}
    \label{tab:citation-results}
    \begin{tabular}{lr}
        \hline
        \textbf{Metric} & \textbf{Result} \\
        \hline
        Citation Precision & \(62.69\%\) \\
        Citation Recall & \(64.80\%\) \\
        Citation F1 & \(62.86\%\) \\
        Citation Exact Match & \(43.20\%\) \\
        \hline
    \end{tabular}
\end{table}

Citation Recall exceeds Citation Precision by a small margin. The system
usually retrieves most gold articles but sometimes adds an article from a
different candidate path or omits a legal basis on the gold path. Citation
Exact Match reaches \(43.20\%\), showing that matching the complete set of
legal bases remains more difficult than retrieving at least one relevant
article.

\subsection{Results by Question Type}

Table~\ref{tab:question-type-results} presents the principal metrics by
question type. R@5, E@5, T1, OR, VA, AF1, and CF1 denote Rule Recall@5,
Event Recall@5, Top-1 Exact Path, Oracle Exact Path, Verification Accuracy,
Answer Token F1, and Citation F1, respectively.

\begin{table*}[t]
    \centering
    \scriptsize
    \caption{Results by question type.}
    \label{tab:question-type-results}
    \begin{tabular}{lrrrrrrrr}
        \hline
        \textbf{Question type} &
        \textbf{N} &
        \textbf{R@5} &
        \textbf{E@5} &
        \textbf{T1} &
        \textbf{OR} &
        \textbf{VA} &
        \textbf{AF1} &
        \textbf{CF1} \\
        \hline
        Branch reasoning &
        \(4\) &
        \(60.00\) &
        \(60.00\) &
        \(0.00\) &
        \(0.00\) &
        \(100.00\) &
        \(69.33\) &
        \(43.57\) \\

        Bridge-convergence fallback &
        \(6\) &
        \(100.00\) &
        \(100.00\) &
        \(83.33\) &
        \(100.00\) &
        \(100.00\) &
        \(65.03\) &
        \(91.67\) \\

        Bridge event &
        \(40\) &
        \(81.25\) &
        \(86.67\) &
        \(50.00\) &
        \(95.00\) &
        \(100.00\) &
        \(75.95\) &
        \(75.42\) \\

        Chain explanation &
        \(35\) &
        \(78.57\) &
        \(82.86\) &
        \(40.00\) &
        \(80.00\) &
        \(100.00\) &
        \(77.22\) &
        \(67.62\) \\

        Convergence reasoning &
        \(10\) &
        \(69.33\) &
        \(82.50\) &
        \(0.00\) &
        \(0.00\) &
        \(100.00\) &
        \(84.32\) &
        \(59.66\) \\

        Forward multihop &
        \(70\) &
        \(69.29\) &
        \(79.52\) &
        \(35.71\) &
        \(87.14\) &
        \(100.00\) &
        \(76.89\) &
        \(68.10\) \\

        Reverse reasoning &
        \(45\) &
        \(46.67\) &
        \(42.22\) &
        \(8.89\) &
        \(26.67\) &
        \(60.00\) &
        \(32.67\) &
        \(36.32\) \\

        Yes/no counterexample &
        \(20\) &
        \(82.50\) &
        \(88.33\) &
        \(55.00\) &
        \(95.00\) &
        \(100.00\) &
        \(89.16\) &
        \(67.50\) \\

        Yes/no positive &
        \(20\) &
        \(67.50\) &
        \(70.00\) &
        \(45.00\) &
        \(70.00\) &
        \(90.00\) &
        \(64.25\) &
        \(62.94\) \\
        \hline
    \end{tabular}
\end{table*}

The bridge event, causal chain explanation, and forward multihop groups
achieve Answer Token F1 scores of \(75.95\%\), \(77.22\%\), and
\(76.89\%\), respectively. Convergence reasoning achieves the highest
Answer Token F1 among the non-counterexample questions, at \(84.32\%\).

Yes/no counterexample achieves an Answer Token F1 of \(89.16\%\) and a
Verification Accuracy of \(100\%\). This result reflects the ability to
distinguish a direct relation from an indirect path. However, this group
does not provide complete ground truth for hard interventions.

Reverse reasoning is the weakest group:

\begin{itemize}
    \item Rule Recall@5: \(46.67\%\);
    \item Event Recall@5: \(42.22\%\);
    \item Top-1 Exact Path: \(8.89\%\);
    \item Oracle Exact Path: \(26.67\%\);
    \item Verification Accuracy: \(60.00\%\);
    \item Answer Token F1: \(32.67\%\);
    \item Citation F1: \(36.32\%\).
\end{itemize}

These results show that reverse-reasoning errors emerge as early as the
retrieval stage and subsequently propagate to path selection, verification,
and answer generation.

Both Top-1 and Oracle Exact Path are \(0\%\) for branch reasoning and
convergence reasoning. Nevertheless, these groups achieve Answer Token F1
scores of \(69.33\%\) and \(84.32\%\), respectively. This discrepancy shows
that a linear path metric is unsuitable for fully evaluating branching or
convergent outputs.

\subsection{Results by Difficulty}

Table~\ref{tab:difficulty-results} presents the results by difficulty level.

\begin{table}[t]
    \centering
    \caption{Results by difficulty level.}
    \label{tab:difficulty-results}
    \begin{tabular}{lrr}
        \hline
        \textbf{Metric} & \textbf{Medium} & \textbf{Hard} \\
        \hline
        Number of questions & \(181\) & \(69\) \\
        Rule Recall@5 & \(67.13\%\) & \(77.29\%\) \\
        Event Recall@5 & \(71.45\%\) & \(83.07\%\) \\
        Top-1 Exact Path & \(34.81\%\) & \(36.23\%\) \\
        Oracle Exact Path & \(72.38\%\) & \(68.12\%\) \\
        Verification Accuracy & \(88.95\%\) & \(100.00\%\) \\
        Answer Token F1 & \(63.90\%\) & \(81.25\%\) \\
        Citation F1 & \(62.02\%\) & \(65.04\%\) \\
        \hline
    \end{tabular}
\end{table}

The hard group achieves higher answer and verification results than the
medium group. This finding should not be interpreted as evidence that
higher difficulty improves system performance. The difficulty label is
correlated with question type: the hard group contains many counterexample,
chain, or structurally templated samples whose gold answers are easier to
match, whereas the medium group contains most reverse-reasoning questions.
Consequently, results by difficulty are affected by data composition and
must be analyzed together with question type.

\subsection{Results on the Counterfactual Subset}

On the \(20\) samples with
\texttt{requires\_counterfactual=true}, the system achieves:

\begin{itemize}
    \item Verification Accuracy: \(100.00\%\);
    \item Answer Token F1: \(89.16\%\);
    \item Citation F1: \(67.50\%\);
    \item Rule Recall@5: \(82.50\%\);
    \item Event Recall@5: \(88.33\%\);
    \item Top-1 Exact Path: \(55.00\%\);
    \item Oracle Exact Path: \(95.00\%\).
\end{itemize}

The gap between Oracle and Top-1 Exact Path on this subset is \(40\)
percentage points. Although the verifier correctly classifies all
direct-edge negative claims, path ranking does not always select the gold
path.

As noted in Section~\ref{sec:data}, these samples primarily test the
rejection of a direct edge. Therefore, the \(100\%\) result is not used to
claim that LegalSCM has fully solved hard-intervention reasoning.

\subsection{Paired Ablation}

Table~\ref{tab:paired-ablation} compares the structural SCM and path ablation
in the paired setup.

\begin{table*}[t]
    \centering
    \caption{Paired comparison of the structural SCM and path ablation.}
    \label{tab:paired-ablation}
    \begin{tabular}{lrrr}
        \hline
        \textbf{Metric} &
        \textbf{Structural SCM} &
        \textbf{Path ablation} &
        \textbf{Difference} \\
        \hline
        Rule Recall@5 &
        \(69.93\%\) & \(69.93\%\) & \(0.00\) \\
        Event Recall@5 &
        \(74.66\%\) & \(74.66\%\) & \(0.00\) \\
        Top-1 Exact Path &
        \(35.20\%\) & \(35.20\%\) & \(0.00\) \\
        Oracle Exact Path &
        \(71.20\%\) & \(71.20\%\) & \(0.00\) \\
        Verification Accuracy &
        \(92.00\%\) & \(92.00\%\) & \(0.00\) \\
        Verification Macro-F1 &
        \(97.73\%\) & \(97.73\%\) & \(0.00\) \\
        Answer Token F1 &
        \(68.69\%\) & \(68.69\%\) & \(0.00\) \\
        Answer ROUGE-L F1 &
        \(67.67\%\) & \(67.67\%\) & \(0.00\) \\
        Citation F1 &
        \(62.86\%\) & \(62.86\%\) & \(0.00\) \\
        Counterfactual Accuracy &
        \(100.00\%\) & \(100.00\%\) & \(0.00\) \\
        \hline
        Mean time/sample &
        \(3.890\) s & \(1.205\) s & \(+2.685\) s \\
        Median time/sample &
        \(2.269\) s & \(1.220\) s & \(+1.049\) s \\
        Batch wall-clock time &
        \(1{,}284.455\) s & \(450.429\) s & \(+834.026\) s \\
        \hline
    \end{tabular}
\end{table*}

The two verification engines produce identical quality metrics on the
current benchmark. Because retrieval is performed before Step 4, identical
retrieval and path metrics for the two configurations are expected. More
importantly, the verification, answer, and citation metrics also remain
unchanged.

For runtime, the ratio of mean time per sample is:

\begin{equation}
\frac{
    3.890
}{
    1.205
}
\approx
3.23.
\end{equation}

Thus, the structural SCM is approximately \(3.23\) times slower in terms of
per-sample runtime. For batch wall-clock time, the ratio is:

\begin{equation}
\frac{
    1{,}284.455
}{
    450.429
}
\approx
2.85.
\end{equation}

The difference between these ratios may arise from initialization, caching,
checkpointing, and artifact-writing costs that are not allocated identically
in the per-sample runtime.

The paired result is regarded as a \emph{null quality ablation}: the
structural SCM does not improve any quality metric on the current benchmark.
The structural SCM offers hard-intervention semantics and more detailed
traces, but the current benchmark does not directly evaluate these benefits.

\subsection{Answers to the Research Questions}

\paragraph{RQ1 --- Retrieval quality.}

The system achieves Rule Recall@5 of \(69.93\%\) and Event Recall@5 of
\(74.66\%\). Retrieval provides relatively high coverage but low precision
because it retains many candidate rules.

\paragraph{RQ2 --- Causal path quality.}

Top-1 Exact Path, the deployable path-selection metric, reaches only
\(35.20\%\), whereas the diagnostic Oracle Exact Path reaches \(71.20\%\).
The result identifies path ranking as a larger bottleneck than candidate-path
generation in many cases; the oracle value is diagnostic rather than a
deployable system result.

\paragraph{RQ3 --- Verification quality.}

The query-aware verifier achieves an Accuracy of \(92.00\%\), correctly
classifies all \(20\) direct-edge negative claims, and assigns \(20\) gold
supported samples to uncertain.

\paragraph{RQ4 --- Answer quality.}

Answer Token F1 reaches \(68.69\%\), ROUGE-L F1 reaches \(67.67\%\), and
Citation F1 reaches \(62.86\%\). Reverse reasoning is the lowest-performing
group.

\paragraph{RQ5 --- Effect of the verification engine.}

The structural SCM and path ablation have identical quality metrics, while
the structural SCM is approximately \(3.23\) times slower per sample. The
current benchmark is not sufficiently sensitive to measure an accuracy
benefit from hard-intervention semantics.
